\section{Supplementary Materials}

\subsection{C++ Source Code: Linked List}
The following code snippet demonstrates a basic implementation of a Singly Linked List in C++.

\begin{lstlisting}[language=C++, caption={Basic implementation of a Singly Linked List in C++}]
#include <iostream>

struct Node {
    int data;
    Node* next;
};

class LinkedList {
public:
    LinkedList() : head(nullptr) {}

    void insert(int val) {
        Node* newNode = new Node{val, head};
        head = newNode;
    }

    void display() {
        Node* temp = head;
        while (temp != nullptr) {
            std::cout << temp->data << " -> ";
            temp = temp->next;
        }
        std::cout << "nullptr" << std::endl;
    }

private:
    Node* head;
};

int main() {
    LinkedList list;
    list.insert(30);
    list.insert(20);
    list.insert(10);
    list.display();
    return 0;
}
\end{lstlisting}

\subsection{Database Table Examples}
Below is an example of the database schema for the NPC system.

\begin{table}[H]
\centering
\caption{The NPC personality and context table structure}
\label{tab:db_npc_table}
\begin{tabularx}{\textwidth}{|l|l|X|}
\hline
\textbf{Field Name} & \textbf{Data Type} & \textbf{Description} \\ \hline
\texttt{npc\_id} & INTEGER (PK) & Unique identifier for the NPC. \\ \hline
\texttt{name} & VARCHAR(50) & Name of the NPC. \\ \hline
\texttt{personality} & TEXT & High-level personality trait description. \\ \hline
\texttt{knowledge} & TEXT & Base knowledge and lore known by the NPC. \\ \hline
\texttt{structure\_id} & INTEGER (FK) & ID of the structure the NPC resides in. \\ \hline
\end{tabularx}
\end{table}

\subsection{Project Calendar}
\begin{landscape}
% Below is a LaTeX generated example using pgfgantt to simulate the requirement. 
% You can manually adjust the "progress" value (0 to 100) to create the solid black bars.

\begin{ganttchart}[
    vgrid,
    expand chart=\linewidth, % Expands chart to fill the landscape page
    y unit title=0.5cm,
    title height=1,
    bar/.append style={fill=black}, % Solid black bars for completion
    bar incomplete/.append style={fill=white, draw=black}, % Empty bars for remaining
    progress=today,
    today=6 % Set this to the current time unit
]{1}{20}
    \gantttitle{Project Title)}{20} \\
    % Row 1: The Months (Dates)
    % Syntax: \gantttitle{Label}{Number of Columns it spans}
    \gantttitle{2004/2}{4} 
    \gantttitle{2004/3}{5} 
    \gantttitle{2004/4}{4} 
    \gantttitle{2004/5}{4} 
    \gantttitle{2004/6}{3} \\ % The "\\" moves to the next row
    
    % Row 2: The Weeks
    % We repeat 1,2,3,4 three times to match the months above
    \gantttitlelist{1,2,3,4}{1} 
    \gantttitlelist{1,2,3,4,5}{1} 
    \gantttitlelist{1,2,3,4}{1} 
    \gantttitlelist{1,2,3,4}{1} 
    \gantttitlelist{1,2,3}{1} \\
    
    % Work Package I
    \ganttgroup{I. Work Package I (60\%)}{1}{4} \\
    \ganttbar[progress=100]{a. Task A}{1}{2.5} \\
    \ganttbar[progress=20]{b. Task B}{2}{3} \\
    \ganttbar[progress=0]{c. Task C}{3}{4} \\
    
    % Work Package II
    \ganttgroup{II. Work Package II (0\%)}{4}{8} \\
    \ganttbar[progress=0]{a. Task A}{4}{6} \\
    \ganttbar[progress=0]{b. Task B}{4}{6} \\
    \ganttbar[progress=0]{c. Task C}{6}{8} \\

    % Work Package III
    \ganttgroup{III. Work Package III (20\%)}{8}{12} \\
    \ganttbar[progress=0]{Task...}{8}{12} 
    
\end{ganttchart}
\end{landscape}

\subsection{Project Expenditures}

\noindent \textit{(Fill in the table below. Include all expenses since the project's launch and express them in TL or other foreign currency. Add columns to match the number of work packages if necessary. Also, list all expenditures under the following categories. Compare them to the budget plan given in the proposal and discuss any discrepancies.)}

\vspace{1em}

\begin{table}[H]
\centering
\begin{tabularx}{\textwidth}{|l|X|X|X|X|c|}
\hline
\multicolumn{1}{|c|}{} & \multicolumn{4}{c|}{\textbf{Work Packages}} & \textbf{} \\ \cline{2-5}
\multicolumn{1}{|c|}{} & \centering\textbf{I} & \centering\textbf{II} & \centering\textbf{III} & \centering\textbf{IV} & \textbf{TOTAL} \\ \hline
\textbf{Equipment} & & & & & \\ \hline
\textbf{Consumable goods} & & & & & \\ \hline
\textbf{Publications/Software} & & & & & \\ \hline
\textbf{Transportation} & & & & & \\ \hline
\textbf{Services} & & & & & \\ \hline
\textbf{TOTAL} & & & & & \\ \hline
\end{tabularx}
\end{table}

\vspace{1em}

\noindent \textbf{Equipment:}
\begin{itemize}[label={}]
    \item Description of equipment purchased...
\end{itemize}
\vspace{1em}

\noindent \textbf{Consumable goods:}
\begin{itemize}[label={}]
    \item Details...
\end{itemize}
\vspace{1em}

\noindent \textbf{Publications/Software:}
\begin{itemize}[label={}]
    \item Details...
\end{itemize}
\vspace{1em}

\noindent \textbf{Transportation:}
\begin{itemize}[label={}]
    \item Details...
\end{itemize}
\vspace{1em}

\noindent \textbf{Services:}
\begin{itemize}[label={}]
    \item Details...
\end{itemize}
